\documentclass{beamer}[10]
\usepackage{pgf}
\usepackage{beamerthemesplit}

% Basado en: http://www.tjansson.dk/?p=82

\usepackage[utf8x]{inputenc}
\usepackage{moreverb}
\usepackage{hyperref}
\usepackage{xcolor}
\usepackage{listings}
\definecolor{mygreen}{rgb}{0,0.6,0}
\definecolor{mygray}{rgb}{0.5,0.5,0.5}
\definecolor{mymauve}{rgb}{0.58,0,0.82}
\definecolor{mygray}{rgb}{0.67, 0.73, 0.78}
\lstset{
	language=Python,
	basicstyle=\footnotesize,
	columns=fullflexible, % http://tex.stackexchange.com/questions/4911/phantom-spaces-in-listings-pdf
	upquote=true,
	breaklines=true,
	commentstyle=\color{mygreen},
	keywordstyle=\color{blue},
	stringstyle=\color{mymauve},
    extendedchars=true, 
    inputencoding=utf8x,
    literate={á}{{\'a}}1
             {é}{{\'e}}1
             {í}{{\'i}}1
             {ó}{{\'o}}1
             {ú}{{\'u}}1
             {ñ}{{\~n}}1
}

\definecolor{kugreen}{RGB}{149,122,143} %{50,93,61}
\definecolor{kugreenlys}{RGB}{132,158,139}
\definecolor{kugreenlyslys}{RGB}{173,190,177}
\definecolor{kugreenlyslyslys}{RGB}{214,223,216}
\setbeamercovered{transparent}
\mode<presentation>
\usetheme[numbers,totalnumber,compress,sidebarshades]{PaloAlto}

% Con número de página, o con footbar
%\setbeamertemplate{footline}[frame number]
\setbeamertemplate{footline}[frame]

  \usecolortheme[named=kugreen]{structure}
  \useinnertheme{circles}
  \usefonttheme[onlymath]{serif}
  \setbeamercovered{transparent}
  \setbeamertemplate{blocks}[rounded][shadow=true]
  \logo{\includegraphics[width=1.5cm]{imagenes/linti.pdf}}

% Ocultar barra de navegación
\beamertemplatenavigationsymbolsempty

\title{Programando con Python y Robots}
%\subtitle{la seguridad, la adecuada instalación y la libertad de elección}
\author{Fernando López}
\institute{LINTI \\ Facultad de Informática \\ Universidad Nacional de la Plata
}
\date{29 de Agosto de 2013}


\begin{document}

\frame{\titlepage \vspace{-0.5cm}

}


\frame
{
\frametitle{Indice}
\tableofcontents[pausesection]
}

\begin{frame}{Los inicios del proyectos}
\begin{itemize}
	\item Basado en los cursos de 
\end{itemize}
\end{frame}
\begin{frame}{Pasantías}
\end{frame}
\begin{frame}{El proyecto en las escuelas}
\end{frame}
\begin{frame}[fragile]{Software necesario}
	Para trabajar va a ser necesario:
	\begin{itemize}[<+->]
		\item Alguna distribución GNU/Linux (de preferencia Lihuen o Debian)
		\item Python 2.6/2.7
		\item Algún intérprete como pycrust o idle (opcional)
		\item Algún editor de texto con resaltado de sintaxis (como Geany/Gedit)
		\item La API para controlar el robot:
			\begin{itemize}
				\item En Lihuen:
\verb+apt-get install robot+
				\item En cualquier distro:
\verb+pip install duinobot+
				\item Para desarrollo:
{\footnotesize
\begin{verbatim}
git clone git@github.com:Robots-Linti/duinobot.git
cd duinobot
python setup.py install
\end{verbatim}
}
			\end{itemize}
	\end{itemize}
\end{frame}

\begin{frame}[fragile]{Primeros pasos}
	\begin{itemize}[<+->]
		\item Descargar el manual:\\
		{\footnotesize
		\url{http://robots.linti.unlp.edu.ar/material\_disponible}}
		\item Abrir PyCrust y conectarnos al robot:
\begin{block}{}
\begin{lstlisting}
from duinobot import *
board = Board()
robot = Robot(board, 0)
\end{lstlisting}
\end{block}
		\item Por defecto los robots vienen con id 0, hay que inicializarlos de a uno:
\begin{block}{}
\begin{lstlisting}
robot.setId(4)
robot = Robot(board, 4)
\end{lstlisting}
\end{block}
	\end{itemize}
\end{frame}
\begin{frame}[fragile]{Movimientos básicos y beep}
	\begin{itemize}
		\item \verb+robot.forward(velocidad, tiempo)+
		\item \verb+robot.backward(velocidad, tiempo)+
		\item \verb+robot.turnLeft(velocidad, tiempo)+
		\item \verb+robot.turnRight(velocidad, tiempo)+
		\item \verb+robot.beep(frecuencia, tiempo)+
	\end{itemize}
	Con tiempo en segundos...
\end{frame}
\begin{frame}[fragile]{Con Joystick}
	\begin{itemize}[<+->]
		\item Principalmente usado para motivar
		\item Acciones predefinidas
\begin{block}{}
\begin{lstlisting}
joy = Joystick(robot, nroJoystick)
joy.play()
\end{lstlisting}
\end{block}
		\item El número de Joystick va de 0 en adelante
		\item La API aún no tiene una forma directa de programar los botones
		\item Los acelerómetros de las Notebooks a veces son listados por el PyGame como Joysticks
		\item joy.play() es un loop que termina con el botón ``start'' del joystick
	\end{itemize}
\end{frame}

\begin{frame}[fragile]{Métodos y funciones útiles}
	\begin{itemize}
		\item board.report() $\leftarrow$ lista de robots encendidos
		\item joysticks() $\leftarrow$ lista de Joysticks disponibles
		\item robot.senses() $\leftarrow$ panel con los valores de los sensores~\footnote{es un poco propenso a generar errores aún, en la versión 0.12 debería ser más seguro}
	\end{itemize}\pause
	\begin{block}{Ejemplo: Robot sorpresa}
		\begin{lstlisting}
from random import random
from duinobot import *
board = Board()
robots = board.report()
eleccion = int(random() * len(robots))
robot = Robot(board, robots[eleccion])
robot.forward(50, 2)
		\end{lstlisting}
	\end{block}
\end{frame}
\begin{frame}
	\begin{center}
		\only<1>{\huge ¿Cómo harían para que el robot vaya haciendo beeps a mientras avanza?}
		\only<2>{\huge ¿Sin detenerse en cada beep?}
	\end{center}
\end{frame}
\begin{frame}[fragile]{Beeps mientras el robot se mueve}
	\begin{itemize}[<+->]
		\item Imposible con robot.forward(velocidad, tiempo)
		\item Decimos que el anterior es bloqueante ya que bloquea el interprete hasta terminar
		\item robot.forward(velocidad) es la versión no bloqueante
		\item Cuando usamos la anterior hay que recordar usar robot.stop() al final
		\item Todos los métodos de movimiento se pueden invocar de forma bloqueante y no bloqueante
		\item También el beep() (se detiene con beep(0))
	\end{itemize}\pause
	\begin{block}{}
		\begin{lstlisting}
robot.forward(50)
for i in range(5):
  robot.beep(i * 100, 1)

robot.stop()
		\end{lstlisting}
	\end{block}
\end{frame}

\begin{frame}
	\begin{itemize}
		\item Duinobot $\rightarrow$ Lihuen/GitHub/Pip
		\item Firmata $\rightarrow$ GitHub/Preinstalado
	\end{itemize}
	\begin{itemize}
		\item Duinobot es la API en Python con la que estaremos en contacto.
		\item Duinobot incluye una copia modificada de PyFirmata encargada de la comunicación.
		\item Interfaz serial.
		\item El robot ejecuta una copia modificada de Firmata.
	\end{itemize}
\end{frame}
\begin{frame}{Primeros pasos}
	\begin{itemize}
		\item Libro.
		\item Intérprete.
		\item Conexión
		\item Beep.
		\item Movimientos.
		\item Movimientos simples.
		\item Figuras.
		\item Sensores y condicionales.
		\item Control remoto con raw\_input.
	\end{itemize}
\end{frame}
\begin{frame}{Algoritmos simples}
	\begin{itemize}
		\item Detenerse en obstáculo.
		\item Esquivar.
		\item Escapar.
		\item Mantenerse en el perímetro.
		\item Seguir línea.
		\item Beep proporcional a la distancia.
		\item Beep inversamente proporcional.
	\end{itemize}
\end{frame}
\frame
{
\frametitle{Direcciones}

\begin{block}{Direcciones de contacto}
\begin{center}
flopez@linti.unlp.edu.ar \\
soportelihuen@linti.unlp.edu.ar
\end{center}
\end{block}

}


%~ \frame
%~ {
%~ \frametitle{Entorno de pre-boot del BIOS (3)}
	%~ Hoy en día sería deseable:
	%~ \begin{itemize}
		%~ \item Acceso al stack TCP, por ejemplo para gestores de 
		%~ arranque por red (en lugar de depender de la option ROM con 
		%~ soporte PXE que pudiera tener la NIC).
		%~ \item Uso de dispositivos como un mouse.
		%~ \item La posibilidad de ampliar o mejorar el soporte de 
		%~ hardware sin option ROMs.
		%~ \item Acceso a direccionamiento de 32 o 64 bits.
	%~ \end{itemize}
%~ }
%~ 
%~ \frame{
%~ 
%~ Efteråret 2004
%~ \subsection{Valg af projekt og vejleder}
%~ \frametitle{Tilfældighed}
%~ 
%~ \begin{block}{Kryptering med kaotiske kredsløb}
%~ \begin{columns}
%~ \column{.3\textwidth} \hspace{0.5cm}
%~ \includegraphics[width=0.7\textwidth]{imagenes/linti-f} 
%~ \column{.7\textwidth}
%~ \textit{Mogens Høgh Jensen}, NBI
%~ \end{columns}
%~ \end{block}
%~ 
%~ \begin{block}{Sandriller}
%~ \begin{columns}
%~ \column{.3\textwidth} \hspace{0.5cm}
%~ \includegraphics[width=0.7\textwidth]{imagenes/linti-f} 
%~ \column{.7\textwidth}
%~ \textit{Clive Ellegaard}, NBI og \emph{Tomas Bohr}, DTU
%~ \end{columns}
%~ \end{block}
%~ }





\end{document}
